% Deep Learning — see issue #24.
\subsection{Deep Learning}
\label{subsec:deep_learning}

\emph{Deep learning} is the branch of machine learning that uses neural networks with many hidden layers, called \emph{deep neural networks}. Its defining feature is \emph{representation learning}: instead of relying on features designed by hand, such as the volume of a brain region measured beforehand, the network learns the features directly from the raw data. Each layer transforms the output of the previous one into a slightly more abstract representation, so that simple patterns found in the first layers are combined into increasingly complex ones in the deeper layers \parencite{alzubaidi2021review}.

Although the main ideas are decades old, deep learning grew rapidly with the increase in available data and in computing hardware \parencite{alzubaidi2021review}. In particular, graphics processing units (GPUs) sped up training by more than an order of magnitude, and between 2010 and 2012 they allowed deep networks to set records and win image recognition contests \parencite{schmidhuber2015deep}. Improvements in how networks are built also contributed, such as the ReLU activation, which reduces the \emph{vanishing gradient} problem, where gradients shrink as they are propagated back through many layers which makes the first layers hard to train \parencite{alzubaidi2021review}. Deep learning, and convolutional networks in particular, have since become dominant in computer vision and are increasingly applied in radiology \parencite{yamashita2018convolutional}, including the classification of Alzheimer's disease from brain MRI \parencite{wen2020convolutional}.

The parameters of a deep network are trained with the same principle as in \cref{subsec:neural_networks}, but plain SGD is often replaced by adaptive optimizers. A widely used one is \emph{Adam}, which keeps a running average of the gradients and of their squares for every parameter and uses them to give each parameter its own adaptive step size. It is computationally efficient, and its own settings typically require little tuning \parencite{kingma2015adam}. Other settings that are not learned, such as the learning rate, the batch size or the number of layers, are called \emph{hyperparameters} and are chosen by the user \parencite{yamashita2018convolutional}.

The goal of training is not to fit the training data but to \emph{generalize}, that is, to predict correctly on data the network has never seen. Because deep networks have a large number of parameters, they can learn patterns and noise that are specific to the training set, a problem known as \emph{overfitting}. To detect it, the data are split into three disjoint sets: a \emph{training set} used to update the parameters, a \emph{validation set} used to choose hyperparameters and monitor overfitting, and a \emph{test set} used only once to report the final performance \parencite{yamashita2018convolutional}. In neuroimaging, where each subject may contribute several scans or slices, the split must be done at the subject level, otherwise images of the same brain end up in more than one set. \textcite{wen2020convolutional} found that about half of the Alzheimer's disease classification studies they reviewed may have suffered from this kind of \emph{data leakage}, which biases the reported performance upward.

The methods used to reduce overfitting are known as \emph{regularization}. The most common ones are summarized in \cref{tab:regularization} \parencite{yamashita2018convolutional}. Among them, \emph{dropout} randomly deactivates a fraction of the neurons at each training step, so the network cannot depend on any single neuron, which reduces overfitting \parencite{srivastava2014dropout}.

% Regularization techniques — used in sections/theoretical_framework/deep_learning.tex.
\begin{table}[t]
  \centering
  \caption{Common regularization techniques.}
  \label{tab:regularization}
  \small
  \begin{tabular}{@{}p{0.32\columnwidth}p{0.6\columnwidth}@{}}
    \toprule
    Technique & Key idea \\
    \midrule
    Weight decay      & Adds a penalty on the size of the weights to the loss. \\
    Dropout           & Randomly sets the output of a fraction of the neurons to zero at each training step. \\
    Data augmentation & Enlarges the training set with modified copies of the images, such as small rotations, shifts or flips. \\
    Batch normalization & Normalizes the inputs of each layer, which improves the flow of gradients and allows higher learning rates. \\
    \bottomrule
  \end{tabular}
\end{table}


Overfitting is a particular concern in medical imaging, where large, well-labeled datasets are rarely available. A common strategy is \emph{transfer learning}: a network is first trained on a large dataset, often of natural images such as ImageNet, and its parameters are then used as the starting point for the medical task, where they are adjusted with the smaller dataset in a step called \emph{fine-tuning} \parencite{yamashita2018convolutional}. When the input is an image, the type of deep network most widely used for these tasks is the convolutional neural network, described in \cref{subsec:cnns}.

